% !TeX root = t3_3000.tex

\begin{figure*}[!b]
    \centering
    \includesvg[width=\textwidth]{Images/Ringbuffer.svg}
    \caption{Das Diagramm zeigt, wie die Videodaten von der Videoquelle (Kamera) zum Endpunkt (Server) verlaufen könnten: Es werden die Videodaten der Kamera in dem rohen Format YUYV an den HB10 weitergegeben, 
             welcher immer etwa 3 Sekunden dieser Daten in einem Ringbuffer speichert. Bei einem Treffer mit dem Dart (Trigger) werden die Daten zunächst verarbeitet und anschließend komprimiert.
             Schließlich sendet der HB10 die komprimierten Videos an den Server. YUYV wurde als initiales Videoformat gewählt, da somit auf die Dekodierung der Videodaten innerhalb des HB10 verzichtet werden kann.}
    \label{fig:ringbuffer}
\end{figure*}

\section{Video- Aufnahme und Komprimierung} \label{sec:1}

Die Realisierung einer stabilen Videoaufnahme für den HB10 erfordert eine performante Interaktion zwischen Hardware und Software. Da die hohen Datenraten unkomprimierter Videostreams eine stabile und starke Internetverbindung voraussetzen, 
was nicht bei jedem Aufsteller eines HB10 der Fall ist, wird eine Komprimierung der Daten benötigt.\newline

Zur Minimierung der Prozessorlast und des Arbeitsspeicherverbrauchs bieten sich innerhalb der Linux-Umgebung hardwaregestützte Kodierungsverfahren an. Diese verlagern die rechenintensiven Kompressionsalgorithmen vom Hauptprozessor auf 
dedizierte Hardwareeinheiten \parencite{_2026_ffmpeg}. Dies ist insbesondere im Kontext der im HB10 eingebauten APU relevant, 
da diese über spezielle Funktionen zur beschleunigten Videoverarbeitung verfügt \parencite{_r1000, gpuopenlibrariessdks_2025_amf_video_encode_api}. Des Weiteren, ist es auch nicht nötig, den gesamten Spielablauf aufzunehmen. 
Es ist ausreichend, Videoausschnitte der einzelnen Würfe zu speichern. Im Folgenden werden die technischen Grundlagen des Kamerazugriffs sowie die Auswahl geeigneter Kompressionsverfahren unter Berücksichtigung der Hardwarebeschleunigung erläutert.

\subsection{Video-Komprimierungsalgorithmen} \label{sec:compression}

In der Videoverarbeitung wird grundlegend zwischen Intra-Frame- und Inter-Frame-Komprimierung unterschieden. Bei der Intra-Frame-Komprimierung wird jedes Bild einzeln bearbeitet. Hierbei lassen sich Daten einsparen, 
indem Informationen innerhalb eines Bildes zusammengefasst oder Details entfernt werden, die das menschliche Auge kaum wahrnehmen kann. Ein Beispiel hierfür ist der \textit{MJPEG}-Standard, 
der das \textit{YCbCr}-Farbmodell verwendet. Dabei wird die Helligkeit detailliert gespeichert, während die Farbinformationen reduziert werden \parencite{pearson_an}. \newline

Zusätzlich kann die Inter-Frame-Komprimierung genutzt werden, um die zeitlichen Ähnlichkeiten zwischen aufeinanderfolgenden Bildern zu verwenden. Anstatt jedes Bild vollständig zu speichern,
werden oft nur die Unterschiede zum vorherigen Bild aufgezeichnet \parencite{putra_2023_interframe}. Der Standard \textit{H.264} setzt dieses Prinzip durch \textit{I-Frames} und \textit{P-Frames} um. 
Während ein I-Frame das gesamte Bild enthält, speichern die folgenden P-Frames lediglich die Veränderungen zum vorangegangenen Bild \parencite{kalva_the}. Dieses
Vorgehen reduziert die Datenmenge sowie die benötigte Bandbreite erheblich.

\subsection{Hardwaregestützte Komprimierung} \label{sec:hwComp}

Das HB10 besitzt einen AMD Ryzen Embedded R1000 Prozessor als APU und ist ein Linux-basiertes System. Die vorgesehene Kamera (Modell: ELP-USB3MP01H-L29) unterstützt die Video Formate MJPEG \textit{YUYV} und H.264 in verschiedenen Auflösungen und Wiederholraten,
in Abhängigkeit der Formate.\newline

Da die Videofunktion im Hintergrund des eigentlichen Spiels ablaufen soll, darf die Komprimierung nur einen kleinen Teil der CPU-Leistung beanspruchen. Dazu kann man den Grafikprozessors verwenden, da dieser deutlich effizienter bei der Durchführung von Video-Komprimierung ist. 
Der Grafikprozessor des HB10 unterstützt die \textit{Video Core Next} (VCN) Engine von AMD  \parencite{_2025_amd,_r1000}.\newline

Die in der APU integrierte VCN-Einheit stellt dedizierte Hardware-Encoder für Formate wie H.264 bereit, wodurch die Rechenlast von den CPU-Kernen auf spezialisierte Hardwarekomponenten verlagert 
wird \parencite{_r1000, _2025_amd}. Für die Umsetzung unter Linux kann die \textit{Video Acceleration API} (VA-API) als Schnittstelle genutzt werden \parencite{_2026_ffmpeg}. Diese erlaubt es, 
die Hardwarebeschleunigung der APU direkt anzusprechen, um Videoströme mit minimalem Ressourcenverbrauch zu kodieren.

\subsection{Kamera als Videoquelle} \label{sec:camera}

Bisher gab es keine Notwendigkeit für eine Kamera am HB10, deshalb besteht noch keine Schnittstelle für den Zugriff auf diese.
Der Zugriff erfolgt unter Linux standardmäßig über die \textit{Video for Linux} (V4L) Schnittstelle des Kernels \parencite{_2020_part}. Diese API ermöglicht eine direkte Kommunikation mit dem Videogerät, 
erfordert jedoch eine komplexe Handhabung der Puffer und Formate. Um die Implementierung effizienter zu gestalten, könnte man auf das Framework \textit{FFmpeg} zurückgegriffen \parencite{_2018_ffmpeg}. FFmpeg dient als Abstraktionsebene für V4L und 
vereinfacht den Prozess der Videoaufnahme erheblich \parencite{_documentation}. Des Weiteren kann FFmpeg auch für den Zugriff auf die Codecs des Grafikprozessors verwendet werden.\newline

Da das verwendete Kamera Modell komprimierte Video-Formate unterstützt, besteht die Möglichkeit, diese zu verwenden. In unserem Fall möchten wir die Videos aber weiterverarbeiten, 
deshalb empfiehlt sich die Verwendung unkomprimierter Formate, um den Rechenaufwand für eine Dekomprimierung zu vermeiden. Des Weiteren ist es vorgesehen, die Videos auf dem HB10 in einem Ringbuffer zu Speichern (siehe Abbildung \ref{fig:ringbuffer}). Unter der Verwendung komprimierter Video-Formate ist es aber schwierig die Ringbuffer Größe zu definieren,
da jeder einzelne Frame unterschiedlich groß ist, und bei Formaten die Intra-frame-Kompression verwenden, nicht garantiert werden kann, dass zu jeden Zeitpunkt ein vollständiger I-Frame im Buffer enthalten ist.

Das Kamera-Modell, dass eingebaut werden soll unterstützt als unkomprimiertes Format das YUYV-Pixelformat, dieses kann also für die Verarbeitung innerhalb des HB10 verwendet werden, für die Übertragung an den Server sollte ein komprimiertes Format gewählt werden.

\subsection{Analyse der Kompressionsverfahren MJPEG, H.264 und H.265}

Die Auswahl eines geeigneten Kompressionsverfahrens für dieÜbertragung an den Server, ist entscheidend für die Balance zwischen der Auslastung lokaler Ressourcen und der benötigten Netzwerkbandbreite. Da wir ein unkomprimiertes Videoformat für die Verarbeitung auf dem HB10 verwenden, 
spielen die restlichen Formate der Kamera keine Rolle, sondern entscheidend ist, welche Formate von dem Grafikprozessor kodiert werden können. Der Grafikprozessor des HB10 unterstützt die Komprimierung in die Formate H.264 und H.265.\newline

H.264 und H.265 nutzen die Inter-Frame-Kompression, um Redundanzen über die Zeit hinweg zu reduzieren \parencite{putra_2023_interframe}. H.264 bietet eine weitreichende Kompatibilität und eine 
effiziente Reduktion der Dateigröße\parencite{kalva_the}. Da die Hardware des Typs AMD Ryzen Embedded R1000 dieses Format unterstützt, könnte die Kompression 
mittels VA-APIdirekt auf die GPU ausgelagert werden \parencite{_2026_ffmpeg, gpuopenlibrariessdks_2025_amf_video_encode_api}. Dies schont die CPU-Ressourcen und den Arbeitsspeicher während des Betriebs \parencite{_2025_amd}.\newline

H.265 erreicht bei gleicher visueller Qualität eine noch stärkere Kompression als H.264, was die benötigte Bandbreite für die Übertragung an einen Server weiter minimiert \parencite{putra_2023_interframe}. Allerdings ist 
die algorithmische Komplexität von H.265 deutlich höher. Dieses Format ist sinnvoll, wenn die Einsparung von Speicherplatz und Rechenleistung auf dem Server und die Reduktion der Übertragungskosten Vorrang vor der minimalen Latenz auf dem lokalen Gerät haben.\newline

Für die Übertragung an den Server bietet sich H.265 an, da es die benötigte Bandbreite im Vergleich zu den anderen Formaten am besten reduziert.

\vspace{1em}